% ---------------------------------------------------------------------------
% Kapitel 8: Zusammenfassung und Ausblick
% Wortgleich uebernommen aus der Confluence-Seite
% "Kapitel 8: Zusammenfassung und Ausblick (Teilentwurf)"
% (Space BEH, Seiten-ID 58097676), Stand 04.10.2026.
% ---------------------------------------------------------------------------

\chapter{Zusammenfassung und Ausblick}
\label{kap:zusammenfassung}

\section{Zusammenfassung}
\label{sec:gesamtzusammenfassung}

Die vorliegende Arbeit hat ein institutionelles E-Portfolio-Portal für die Hochschule Bremen entworfen und gebaut. Der Weg führte von der Forschungsliteratur (Kapitel 3) und der heuristischen Analyse einer Referenzplattform (Abschnitt 3.4) zu einem Katalog aus 36 funktionalen und 24 nicht-funktionalen Anforderungen (Kapitel 2), von dort zu Konzeption und Realisierung (Kapitel 4 und 5). Die Realisierung diente als Werkzeug, um zu prüfen, ob die abgeleiteten Anforderungen einer Umsetzung standhalten. Geprüft wurde mit vier Methoden und elf Personen (Kapitel 6), an Maßstäben, die vor dem Test feststanden (Kapitel 7).

Von den funktionalen Anforderungen sind 30 erfüllt und sechs teilweise, von den nicht-funktionalen zehn erfüllt und elf teilweise; drei blieben ungeprüft. Der SUS-Mittelwert liegt mit 87,3 um 19,3 Punkte über dem vorab gesetzten Zielwert von 68 (Lewis \& Sauro, 2018). Dem hohen Urteil stehen 30 Usability-Probleme gegenüber, sieben davon schwerwiegend.

\subsection{Antwort auf die erste Teilfrage}

Die erste Teilfrage lautet: „Welche Anforderungen an Struktur, Inhalt und Workflow lassen sich aus der bestehenden Forschungsliteratur zu E-Portfolios ableiten?“ Zur Struktur liefert die Literatur zwei Anforderungen. Lorenzo und Ittelson (2005) bestimmen ein E-Portfolio als Sammlung, die auch eine Institution repräsentieren kann, und benennen die Leserschaft veröffentlichter Unterlagen einer Einrichtung als offen bis hin zu jeder Person mit Internetzugang. Daraus folgen der Zugang ohne Anmeldung und die Gewichtung der externen Nutzenden als Hauptzielgruppe; die Übertragung der offenen Leserschaft auf studentische Arbeiten ist dabei ein Schritt der vorliegenden Arbeit (Abschnitt 3.1). Ordnet man die zehn Komponenten von Mahasneh (2020) nach ihrem Bezugspunkt, was die Quelle selbst nicht tut, ergeben sich eine personenbezogene und eine leistungsbezogene Ebene und damit die Trennung von Profilseite und Projekt. Zum Inhalt liefert dasselbe Modell das Raster: Sieben seiner zehn Komponenten deckt das Portal ab, die drei an eine Lehrveranstaltung gebundenen lässt es mit der Abgrenzung gegenüber dem Lernmanagementsystem und der eigenen Einordnung als Präsentationsportfolio bewusst aus (Tabelle 7.5). Zum Workflow trägt die Bestimmung des institutionellen Portfolios als fokussierte Auswahl bei (Kahn, 2001, zitiert nach Lorenzo \& Ittelson, 2005); die vorliegende Arbeit folgert daraus einen Auswahlschritt vor der Veröffentlichung, also eine redaktionelle Freigabe. Quer zu allen dreien steht die Akzeptanzforschung: Davis (1989) liest die Benutzerfreundlichkeit vorläufig als Vorbedingung der wahrgenommenen Nützlichkeit, und Shroff et al. (2011) stützen diesen Pfad an einem E-Portfolio-System. Die vorliegende Arbeit folgert daraus, dass Gebrauchstauglichkeit eine Frage des Entwurfs ist und nicht der Nacharbeit (Abschnitt 3.2).

Von den neun funktionalen Anforderungen, die Tabelle 2.1 auf die E-Portfolio-Literatur zurückführt, sind acht erfüllt; der neunten (FA-16) fehlt allein der Studiengang der Person, ein Glied aus der Ordnung der Hochschule. Nicht aus der Literatur folgen dagegen die Aufteilung des Freigabewegs in fünf Zustände, die Ableitung der Rolle aus der Hochschuladresse, die Ordnungsstruktur der Hochschule und die kuratierte Startseite. Sie stammen aus dem Anwendungsfall.

\subsection{Antwort auf die zweite Teilfrage}

Die zweite Teilfrage lautet: „Welche funktionalen und gestalterischen Merkmale fördern die Auffindbarkeit und Benutzerfreundlichkeit eines institutionellen Portfolio-Portals?“ Abschnitt 7.2 hat sie aus den Testdaten beantwortet. Die Auffindbarkeit fördern sichtbar angebotene, verbindbare Ordnungsmerkmale. Alle elf Personen, auch die beiden ohne Hochschulbezug, schränkten drei Merkmale zugleich ein, im Median in 59 Sekunden. Hinzu kommen die einzeln rücknehmbare Einschränkung und der zweite Zugang über die Person, den ebenfalls alle elf fanden. Die Benutzerfreundlichkeit fördern kurze Wege für wiederkehrende Vorgänge und ein Editor, der ohne Einweisung bedienbar ist: Einreichen und Zurückziehen dauerten im Median 20 und 25 Sekunden, und die Aufgabe mit mehreren Bausteinarten lösten alle fünf Studierenden in höchstens 3:11 Minuten. Der SUS-Mittelwert liegt in der Note A+ der Skala, die Lewis und Sauro (2018) wiedergeben.

Nicht bestätigt hat sich die Annahme, Trefferzahl und Zählmarken zeigten den Zustand der Suche ausreichend; die wirkende Einschränkung muss bei den Treffern selbst stehen. Neu ist eine Bedingung, unter der die genannten Merkmale wirken: Handlungen und Zustände müssen beschriftet sichtbar sein (H6), und dasselbe muss überall gleich heißen (H4); dort liegen die beiden Schwerpunkte der Problemliste. Gemessen an den Heuristiken nach Nielsen (2024) schließt das Portal gegenüber der Referenzplattform die Lücke bei H7 im öffentlichen Bereich und bleibt bei H2 und H4 zurück, weil es mehr Funktionen zeigt, sie aber weniger einheitlich benennt.

\subsection{Antwort auf die dritte Teilfrage}

Die dritte Teilfrage lautet: „Wie lässt sich ein blockbasiertes Inhaltsmodell mit integriertem Redaktionsworkflow realisieren, das Studierenden eine flexible und medienreiche Projektpräsentation ermöglicht?“ Die Realisierung bestätigt vier Entscheidungen des Entwurfs. Erstens ist das Inhaltsmodell an das Dokumentmodell des Editors gebunden: Ein Baustein entspricht einem Knoten der obersten Ebene, und zwischen Editor und Speicherung genügt eine schmale Zuordnung (Abschnitt 5.2). Zweitens wird der Inhalt eines Bausteins unverändert gespeichert, seine Art aber zusätzlich als abfragbares Merkmal geführt, weil sich daran Regeln und Auswertungen wartbar formulieren lassen. Drittens liegen Dateien außerhalb des Dokuments, ein Baustein verweist nur auf sie, und ausgeliefert werden sie über einen eigenen Weg, weil ihre Sichtbarkeit am Zustand des Projekts hängt. Viertens ist die Bearbeitbarkeit an diesen Zustand gebunden und auf dem Server durchgesetzt, auch für das Hochladen und Entfernen von Dateien. Darin besteht die Integration des Workflows in das Inhaltsmodell. Allein die Administration kann Dateien in jedem Zustand entfernen, was Abschnitt 7.3 als Abweichung von FA-42 führt. Zwei Schranken, die der Entwurf vorsah, hat erst die Umsetzung auf dem Server verbindlich gemacht: Eingereicht wird nur mit zugewiesener Redaktion, erneut eingereicht erst ohne offene Änderungsaufforderung.

Für mehrere Bearbeitende verlangt das automatische Zwischenspeichern ein Versionsfeld, das einen Konflikt meldet, statt fremde Arbeit still zu überschreiben; eine Bearbeitungssperre macht solche Konflikte selten. Die Flexibilität trägt ein Katalog aus fünfzehn Bausteinarten, dessen Bandbreite dem Komponentenmodell von Mahasneh (2020) folgt; Aufgabenlisten fehlen bewusst, weil die vorliegende Arbeit das Portal als Präsentationsportfolio einordnet (Abschnitt 3.1). Dass Flexibilität nicht immer mit Benutzerfreundlichkeit einhergeht, hält Davis (1989) fest, und der Nutzertest zeigt es an einer Stelle: Die Aufgabe mit mehreren Bausteinarten gelang allen fünf Studierenden, doch eine Person fand zwei der Arten nicht rasch (P4). Hinzu kommt ein Fehler der Umsetzung: Der im Editor angekündigte Schrägstrich öffnet das Auswahlmenü nicht (P30).

\subsection{Antwort auf die Forschungsfrage}

Die Forschungsfrage lautet: „Welche funktionalen und gestalterischen Anforderungen muss ein hochschulisches E-Portfolio-Portal erfüllen, um studentische Projekte strukturiert zu präsentieren, auffindbar zu machen und von den beteiligten Nutzergruppen akzeptiert zu werden?“

Strukturiert präsentiert ein Portal, das zwischen Projekt und Person trennt und ein Projekt aus frei anordenbaren Bausteinen aufbaut; jedes Projekt durchläuft vor der Veröffentlichung eine redaktionelle Freigabe. Auffindbar macht das Portal seinen Bestand, wenn es die Ordnung der Hochschule als verbindbare und einzeln rücknehmbare Merkmale anbietet, die wirkende Einschränkung bei den Treffern zeigt und einen Zugang über die Person eröffnet. Akzeptiert wird es, wenn sein Nutzen erkennbar ist, eigene Arbeit also für Externe sichtbar wird, und wenn Zustände wie Handlungen beschriftet sichtbar sind. Bei Nützlichkeit und Benutzerfreundlichkeit liegen alle elf Befragten über dem Orientierungspunkt von 4,00, bei der Nutzungsabsicht zehn darüber und eine Person genau darauf; beschriftet sichtbar sind die Handlungen dagegen erst teilweise (P6, P7). Wie bei Davis (1989) hängt die Nutzungsabsicht stärker mit der wahrgenommenen Nützlichkeit zusammen als mit der Benutzerfreundlichkeit (explorative Rangkorrelation 0,84 gegenüber 0,41). Den Pfad von der Benutzerfreundlichkeit zur Nützlichkeit, den Shroff et al. (2011) stützen, zeigt diese kleine Stichprobe dagegen nicht.

Unter dem Vorbehalt der folgenden Grenzen erscheinen vier Anforderungen übertragbar, die an Aussagen der Literatur und an die Heuristiken nach Nielsen (2024) anschließen: der Zugang ohne Anmeldung, die Freigabe vor der Veröffentlichung, die Trennung von Projekt und Person sowie die sichtbare Beschriftung von Zustand und Handlung. An die Hochschule Bremen gebunden bleiben die Ausgestaltung des Freigabewegs, die Ableitung der Rolle aus der Hochschuladresse, die kuratierte Startseite und die Ordnung aus fünf Fakultäten mit 75 Studiengängen (Abschnitt 2.1), entlang der gefiltert wird.

\subsection{Grenzen der Arbeit}

Die empirische Grundlage ist schmal. Elf Personen nahmen an einer einzigen Testrunde teil. Die Blöcke der Redaktion und der Administration mit je zwei Personen und die beiden externen Nutzenden liegen unter den drei Personen, die Nielsen (2000) je Nutzergruppe mindestens verlangt. Studierende bearbeiteten beide Rollenblöcke, Lehrende nahmen nicht teil, und die Aufgabenkarten erklärten dort den Ablauf; diese Blöcke zeigen deshalb, ob die Oberfläche einen bekannten Ablauf trägt, nicht ob er ohne Anleitung gefunden wird. Unter den Externen war niemand vor der Studienwahl. Die heuristische Evaluation stammt von zwei statt drei bis fünf Begutachtenden (Nielsen \& Molich, 1990), und beide haben das System selbst gebaut; eine Verzerrung zugunsten des eigenen Systems ist wahrscheinlicher als die umgekehrte. Beide Bögen sind übersetzt, aber nicht eigens validiert; der TAM-Bogen erfragt die Nutzungsabsicht nach einer einzigen Sitzung, während Shroff et al. (2011) nach einem Semester Nutzung fragten. Die Teilnehmenden wussten, dass die Testleitenden das System gebaut hatten; die hohen Werte, darunter neun SUS-Werte in der Note A+ und sechs Höchstwerte der Nutzungsabsicht, können daher soziale Erwünschtheit und einen Deckeneffekt enthalten. Aufbauende Aufgaben ließen Reihenfolge- und Lerneffekte zu, die Moderation wechselte nicht wie geplant, und vier Wertungen weichen von den eigenen Regeln ab (Abschnitte 6.1 und 6.3).

Auch die Grundlage der Anforderungen hat Grenzen. Verglichen wurde mit einer einzigen Referenzplattform, die auf Empfehlung und nicht über ein Auswahlverfahren gewählt wurde. Begutachtet wurde sie stichprobenartig und vier Monate früher, allein aus der Sicht ohne Anmeldung und ebenfalls von nur zwei Personen (Nielsen \& Molich, 1990). Eine Befragung künftiger Nutzender vor der Entwicklung fand nicht statt (Abschnitt 2.2). Ungeprüft blieben Antwortzeiten und Tragfähigkeit (NFA-30 bis NFA-32). Die Barrierefreiheit nach W3C (2023) ist nur teilweise belegt: automatisiert, im Ausgangszustand jeder Seite, ohne Vorlesesoftware, ohne Messung der Kontraste von Bedienelementen oder von Text über Bildern und nicht am getesteten Stand protokolliert. Die Darstellung ist nur in einem Browser und ohne Tabletbreite oder echte Geräte belegt (NFA-40, NFA-41). Schließlich ist das System ein Prototyp, getestet mit einem vorbereiteten Bestand von 48 veröffentlichten Projekten.

\section{Ausblick}
\label{sec:ausblick}

Die wirkende Einschränkung gehört über die Trefferliste (P17), und Handlungen hinter Symbolen und Statusmarken brauchen eine eigene, beschriftete Schaltfläche (P1, P2, P6, P7). Begriffe und Bedienmuster sind zu vereinheitlichen. Zu beheben sind außerdem die zu geringen Kontraste von Text über hellen Bildern und von Feldrändern (W3C, 2023), die überlappenden Personenbilder auf den Projektkarten (NFA-24), der Verlust ungesicherter Angaben beim Verlassen eines Formulars (P28) und der wirkungslose Schrägstrich im Editor (P30). Noch offen sind die Suche über die Namen Mitwirkender (FA-14), der Studiengang auf der Profilseite (FA-16), eine Übersicht der Studiengänge und Module (FA-18), Untertitel und Transkripte für Audio und Video sowie der Abgleich der Löschbefugnisse der Administration mit FA-42 (Abschnitt 7.3). Danach wäre erneut zu testen, nach Nielsen (2000) in mehreren kleinen Runden mit Überarbeitung dazwischen und mit mindestens drei Personen je Gruppe, in der Redaktion mit Lehrenden.

Zwei Anregungen kamen außerhalb der Erhebung aus der Hochschule selbst. Die erste betrifft die Startseite mit ihrer einzigen kuratierten Top 10. Die fünf Fakultäten der Hochschule (Abschnitt 2.1) teilen sich also zehn Plätze, und keine Regel sichert jeder Fakultät einen davon. Der Erstprüfer regte an, ganz oben auf der Startseite zufällig gewählte Projekte aus jeder Fakultät zu zeigen (A. Teufel, persönliche Mitteilung, 8. September 2026); in den Sitzungen schlug zudem eine Person eine Top 10 nach Fakultäten vor (Abschnitt 6.3). Ein Portal, das nach Lorenzo und Ittelson (2005) eine Institution repräsentiert, sollte deren Breite schon auf der ersten Seite zeigen. Die vorliegende Arbeit schlägt dafür eine kuratierte Top 10 je Fakultät vor. Eine reine Zufallsauswahl aus dem geprüften Bestand gäbe dagegen die kuratorische Entscheidung auf, welche Arbeiten eine Fakultät vorn zeigt (Abschnitt 4.4). Mit der Anregung des Erstprüfers lässt sich das verbinden, indem ganz oben aus jeder Liste wechselnd ein Projekt erscheint.

Die zweite Anregung betrifft eine Gruppe, die Abschnitt 2.1 unter den externen Nutzenden nicht führt: Schülerinnen und Schüler vor der Studienwahl. Ein Gespräch am 25. September 2026 mit zwei Mitarbeitenden der Zentralen Studienberatung, an die der Erstprüfer verwiesen hatte, gehört nicht zur Erhebung aus Kapitel 6 und wurde nicht ausgewertet. Die beiden schätzten ein, dass diese Gruppe die Projekte nicht ohne Weiteres versteht. Kurzbeschreibung und Textbausteine, die im Portal die Komponente der Zusammenfassungen nach Mahasneh (2020) tragen, liegen nur in der Fassung vor, die von den Studierenden selbst stammt. Gebraucht wird eine eigene Ansicht, die ein Projekt in verständlicher Sprache zusammenfasst; ein Sprachmodell, das solche Fassungen erzeugt, wäre ein möglicher Weg dorthin. Zu klären wäre dabei, wie sich eine maschinell erzeugte Fassung zum Freigabeweg aus Abschnitt 4.4 verhält. Da sich Gebrauchstauglichkeit nach Brooke (1996) nur in einem Nutzungskontext bestimmen lässt, verlangt eine neue Gruppe einen eigenen Test, und ob ihr die Ansicht nützt, ließe sich mit den Konstrukten nach Davis (1989) erheben.

Am Inhaltsmodell ist die gleichzeitige Bearbeitung durch mehrere Mitwirkende der nächste Schritt; ein Entwurf dafür lag vor, umgesetzt wurde stattdessen die Bearbeitungssperre aus Abschnitt 4.3. Kommentare an einzelnen Bausteinen würden zudem die Prüfung erleichtern; sie setzen voraus, dass ein Baustein seine Kennung behält, wenn sich seine Stelle im Projekt verschiebt, was heute nicht gesichert ist.

Für einen Produktivbetrieb fehlen einige Vorkehrungen. Die Abmeldung wirkt allein im Browser. Bearbeitungssperre und Begrenzung wiederholter Anmeldeversuche gelten nur für einen einzelnen Server, und das Löschen eines Kontos ist eine Deaktivierung, der die Anonymisierung der Autorschaft noch fehlt (Abschnitt 4.6). Die Eigenentwicklung trägt zudem nach Lorenzo und Ittelson (2005) das Risiko, das Fachpersonal nicht lange genug zu binden. Weil die Entwicklung an eine Abschlussarbeit gebunden war, setzt ein Betrieb voraus, dass eine Stelle der Hochschule Wartung und Weiterentwicklung übernimmt; API-Contract, Dokumentation und automatisierte Tests (Abschnitte 4.1 und 5.3) sind die Grundlage dafür. Offen bleibt schließlich, was nach Tabelle 6.1 keine der vier Methoden beantwortet: ob aus der gemessenen Nutzungsabsicht tatsächlich Nutzung wird. Das klärt erst eine Erhebung im laufenden Betrieb.
